% GENERATED FILE. Edit canonical lessons, then run npm run generate:books.
% canonical-source-hash: fnv1a64:2f7b5dbba1407faa
% canonical-lessons: SA-C297-indhanam, SA-C297-ghatah, SA-C297-dhara, SA-C297-nirvapayati, SA-C297-samyati

\chapter{Putting the Fire Out}
\label{ch:sa-putting-the-fire-out}
\hlchaptermodality{297}

\begin{chapteropening}
\textbf{By the end of this chapter:} \emph{I can say how a fire was put out.}

Complete the last lesson of chapter 297: I can say how a fire was put out.
\end{chapteropening}

\section[\texorpdfstring{\sk{इन्धनम्} (indhanam)}{इन्धनम् (indhanam)}]{\sk{इन्धनम्} (indhanam) — fuel; firewood}

\label{lesson:SA-C297-indhanam}



\begin{quote}
Before the new one: say the Sanskrit for dense, then the Sanskrit for spreads.
\end{quote}

\subsection*{You'll want to know: \sk{इन्धनम्}}
\textbf{\sk{इन्धनम्}} — \emph{indhanam} — \textquotedblleft{}fuel; firewood\textquotedblright{}.

\begin{quote}
\sk{शुष्कम्} \sk{इन्धनं} \sk{शीघ्रं} \sk{ज्वलति।}
\end{quote}

\begin{quote}
\emph{śuṣkam indhanaṃ śīghraṃ jvalati.}
\end{quote}

\begin{quote}
Dry fuel burns fast.
\end{quote}

A neuter noun.

\begin{cousinweb}
\textbf{\sk{इन्धनम्}} (\emph{indhanam}) comes from the root √indh, \textquotedblleft{}to kindle, light\textquotedblright{}; MW gives \emph{indhana} as kindling and as fuel (MW).
\end{cousinweb}

\begin{grammarlens}[title={What you've built}]
One more word for telling what happened.
\end{grammarlens}

\subsection*{Your turn}
\begin{itemize}
\raggedright
  \item \emph{Say it:} \emph{indhanam}
  \item \emph{Say it:} \emph{indhanam}, once more
  \item \emph{Say it:} \emph{śuṣkam indhanaṃ śīghraṃ jvalati.} — Dry fuel burns fast.
  \item \emph{Recall:} say the Sanskrit for dense, then the Sanskrit for spreads, then say \emph{indhanam} again
\end{itemize}

\begin{tcolorbox}[breakable,colback=teal!4,colframe=teal!35!black,title={Before you move on}]
What does \sk{इन्धनम्} mean? (\textquotedblleft{}Fuel; firewood\textquotedblright{}.) Say it once more.
\end{tcolorbox}
\section[\texorpdfstring{\sk{घटः} (ghaṭaḥ)}{घटः (ghaṭaḥ)}]{\sk{घटः} (ghaṭaḥ) — a water jar}

\label{lesson:SA-C297-ghatah}



\begin{quote}
Before the new one: say the Sanskrit for spreads, then the Sanskrit for fuel.
\end{quote}

\subsection*{You'll want to know: \sk{घटः}}
\textbf{\sk{घटः}} — \emph{ghaṭaḥ} — \textquotedblleft{}a water jar\textquotedblright{}.

\begin{quote}
\sk{जनाः} \sk{घटैः} \sk{जलम्} \sk{आनयन्ति।}
\end{quote}

\begin{quote}
\emph{janāḥ ghaṭaiḥ jalam ānayanti.}
\end{quote}

\begin{quote}
People bring water in jars.
\end{quote}

A masculine noun.

\begin{cousinweb}
\textbf{\sk{घटः}} (\emph{ghaṭaḥ}) MW gives \emph{ghaṭa} as a jar or pitcher, a large earthen water-jar, without a derivation, so its origin is uncertain (MW).
\end{cousinweb}

\begin{grammarlens}[title={What you've built}]
One more word for telling what happened.
\end{grammarlens}

\subsection*{Your turn}
\begin{itemize}
\raggedright
  \item \emph{Say it:} \emph{ghaṭaḥ}
  \item \emph{Say it:} \emph{ghaṭaḥ}, once more
  \item \emph{Say it:} \emph{janāḥ ghaṭaiḥ jalam ānayanti.} — People bring water in jars.
  \item \emph{Recall:} say the Sanskrit for spreads, then the Sanskrit for fuel, then say \emph{ghaṭaḥ} again
\end{itemize}

\begin{tcolorbox}[breakable,colback=teal!4,colframe=teal!35!black,title={Before you move on}]
What does \sk{घटः} mean? (\textquotedblleft{}A water jar\textquotedblright{}.) Say it once more.
\end{tcolorbox}
\section[\texorpdfstring{\sk{धारा} (dhārā)}{धारा (dhārā)}]{\sk{धारा} (dhārā) — a stream of water}

\label{lesson:SA-C297-dhara}



\begin{quote}
Before the new one: say the Sanskrit for fuel, then the Sanskrit for a water jar.
\end{quote}

\subsection*{You'll want to know: \sk{धारा}}
\textbf{\sk{धारा}} — \emph{dhārā} — \textquotedblleft{}a stream of water\textquotedblright{}.

\begin{quote}
\sk{जलस्य} \sk{धारा} \sk{अग्नौ} \sk{पतति।}
\end{quote}

\begin{quote}
\emph{jalasya dhārā agnau patati.}
\end{quote}

\begin{quote}
A stream of water falls on the fire.
\end{quote}

A feminine noun.

\begin{cousinweb}
\textbf{\sk{धारा}} (\emph{dhārā}) MW gives \emph{dhārā} as a stream or current of water, without a root in the entry, so its origin is left uncertain here (MW).
\end{cousinweb}

\begin{grammarlens}[title={What you've built}]
One more word for telling what happened.
\end{grammarlens}

\subsection*{Your turn}
\begin{itemize}
\raggedright
  \item \emph{Say it:} \emph{dhārā}
  \item \emph{Say it:} \emph{dhārā}, once more
  \item \emph{Say it:} \emph{jalasya dhārā agnau patati.} — A stream of water falls on the fire.
  \item \emph{Recall:} say the Sanskrit for fuel, then the Sanskrit for a water jar, then say \emph{dhārā} again
\end{itemize}

\begin{tcolorbox}[breakable,colback=teal!4,colframe=teal!35!black,title={Before you move on}]
What does \sk{धारा} mean? (\textquotedblleft{}A stream of water\textquotedblright{}.) Say it once more.
\end{tcolorbox}
\section[\texorpdfstring{\sk{निर्वापयति} (nirvāpayati)}{निर्वापयति (nirvāpayati)}]{\sk{निर्वापयति} (nirvāpayati) — puts out}

\label{lesson:SA-C297-nirvapayati}



\begin{quote}
Before the new one: say the Sanskrit for a water jar, then the Sanskrit for a stream of water.
\end{quote}

\subsection*{You'll want to know: \sk{निर्वापयति}}
\textbf{\sk{निर्वापयति}} — \emph{nirvāpayati} — \textquotedblleft{}puts out\textquotedblright{}.

\begin{quote}
\sk{जनाः} \sk{जलेन} \sk{अग्निं} \sk{निर्वापयन्ति।}
\end{quote}

\begin{quote}
\emph{janāḥ jalena agniṃ nirvāpayanti.}
\end{quote}

\begin{quote}
People put out the fire with water.
\end{quote}

Past (he …): \sk{सः} \sk{जलेन} \sk{अग्निं} \sk{निर्वापितवान्}, \emph{saḥ jalena agniṃ nirvāpitavān}.

\begin{cousinweb}
\textbf{\sk{निर्वापयति}} (\emph{nirvāpayati}) is built from \emph{nir}, \textquotedblleft{}out\textquotedblright{}, and the causative of the root √vā, \textquotedblleft{}to blow\textquotedblright{}; MW gives \emph{nir-vāpayati} as to put out, extinguish (MW).
\end{cousinweb}

\begin{grammarlens}[title={What you've built}]
One more word for telling what happened.
\end{grammarlens}

\subsection*{Your turn}
\begin{itemize}
\raggedright
  \item \emph{Say it:} \emph{nirvāpayati}
  \item \emph{Say it:} \emph{nirvāpayati}, once more
  \item \emph{Say it:} \emph{janāḥ jalena agniṃ nirvāpayanti.} — People put out the fire with water.
  \item \emph{Recall:} say the Sanskrit for a water jar, then the Sanskrit for a stream of water, then say \emph{nirvāpayati} again
\end{itemize}

\begin{tcolorbox}[breakable,colback=teal!4,colframe=teal!35!black,title={Before you move on}]
What does \sk{निर्वापयति} mean? (\textquotedblleft{}Puts out\textquotedblright{}.) Say it once more.
\end{tcolorbox}
\section[\texorpdfstring{\sk{शाम्यति} (śāmyati)}{शाम्यति (śāmyati)}]{\sk{शाम्यति} (śāmyati) — dies down}

\label{lesson:SA-C297-samyati}



\begin{quote}
Before the new one: say the Sanskrit for a stream of water, then the Sanskrit for puts out a fire.
\end{quote}

\subsection*{You'll want to know: \sk{शाम्यति}}
\textbf{\sk{शाम्यति}} — \emph{śāmyati} — \textquotedblleft{}dies down\textquotedblright{}.

\begin{quote}
\sk{शनैः} \sk{अग्निः} \sk{शाम्यति।}
\end{quote}

\begin{quote}
\emph{śanaiḥ agniḥ śāmyati.}
\end{quote}

\begin{quote}
Slowly the fire dies down.
\end{quote}

Past (it …): \sk{अग्निः} \sk{अन्ते} \sk{शान्तवान्}, \emph{agniḥ ante śāntavān}.

\begin{cousinweb}
\textbf{\sk{शाम्यति}} (\emph{śāmyati}) comes from the root √śam, \textquotedblleft{}to be quiet, come to rest\textquotedblright{}; MW gives it as to cease and to be extinguished (MW).
\end{cousinweb}

\begin{grammarlens}[title={What you've built}]
One more word for telling what happened.
\end{grammarlens}

\subsection*{Your turn}
\begin{itemize}
\raggedright
  \item \emph{Say it:} \emph{śāmyati}
  \item \emph{Say it:} \emph{śāmyati}, once more
  \item \emph{Say it:} \emph{śanaiḥ agniḥ śāmyati.} — Slowly the fire dies down.
  \item \emph{Recall:} say the Sanskrit for a stream of water, then the Sanskrit for puts out a fire, then say \emph{śāmyati} again
\end{itemize}

\begin{tcolorbox}[breakable,colback=teal!4,colframe=teal!35!black,title={Before you move on}]
What does \sk{शाम्यति} mean? (\textquotedblleft{}Dies down\textquotedblright{}.) Say it once more.
\end{tcolorbox}
